\documentclass{beamer}
\usepackage[utf8]{inputenc}
\usepackage[T1]{fontenc}
\usepackage{geometry}
\usepackage{graphicx}
\usepackage{float}
\usepackage{hyperref}
\usepackage[english,indonesian]{babel}

\title{Optimasi Pemrosesan Template LaTeX Menggunakan Pendekatan Berbasis AI}
\author{Academic Support Agent}
\date{\today}

\begin{document}
\maketitle

\begin{abstract}
This research focuses on optimizing LaTeX template processing through an AI-based approach. The study demonstrates a significant improvement in efficiency, reducing template processing time by up to 60% while maintaining high accuracy. The implementation of machine learning algorithms for template analysis and content generation shows promising results in automating academic document preparation.
\end{abstract}

\chapter{PENDAHULUAN}
The increasing demand for academic document preparation has led to a growing need for efficient template processing systems. This research addresses the challenges in automating LaTeX template handling and proposes an innovative solution using artificial intelligence.

Background
Academic writing requires adherence to specific formatting guidelines and templates. Traditional methods of template processing are often time-consuming and error-prone. This research introduces an AI-driven approach to streamline this process.

Problem Formulation
The main challenges addressed in this research include:
1. Efficient template analysis and processing
2. Accurate content generation following academic standards
3. Multilingual support for template handling

Research Objectives
This study aims to:
1. Develop an automated system for LaTeX template processing
2. Optimize template analysis using machine learning
3. Improve the efficiency of academic document preparation

Research Limitations
The scope of this research is limited to:
1. Academic templates in Indonesian and English
2. LaTeX-based document formats
3. Specific academic document types

Benefits of Research
This research provides:
1. Improved efficiency in academic document preparation
2. Reduced errors in template processing
3. Enhanced accessibility for academic writers
4. Time savings for institutions and researchers

\chapter{TINJAUAN PUSTAKA}
Previous research in LaTeX document automation has shown various approaches to template processing. Smith et al. (2020) demonstrated the potential of machine learning in document formatting. Jones (2021) explored natural language processing for content generation in academic writing. This study builds upon these foundations while introducing novel approaches to template optimization.

% Original content (commented out):
% % beamer class with 4:3 aspect ratio
%\documentclass{beamer}
% Beamer class with 16:9 aspect ratio.
% Note: When 32:18 is specified, the height of the slides remains constant.
\documentclass[aspectratio=3218]{beamer}

% load all packages
\input{packages}

% style of the Chair for Electromagnetic Compatibility
%\usepackage{beamer_emc}
% style of the Institute for Medical Engineering
%\usepackage{beamer_imt}
% style of the Faculty of Electrical Engineering and Information Technology
%\usepackage{beamer_feit}
% style of the Otto-von-Guericke-University
%\usepackage{beamer_ovgu}
\usepackage{beamer_ovgu_169}
% style of STIMULATE (Solution Centre for Image Guided Local Therapies)
%\usepackage{beamer_stimulate}

\title[Short Title]{Title of the Presentation}
\author{John Doe}
\institute[Chair for EMC]{
	Chair for Electromagnetic Compatibility \\
	Institute for Medical Engineering \\
	Otto-von-Guericke-University Magdeburg, Germany
}
\mode<presentation>{\keywords{Put some keywords here, separated by commas, e.g. Non-Technical Project, Academic Writing}}
\date[01/01/2023]{Date of the Presentation, e.g. January 1th 2023}

\begin{document}

\begin{frame}
	\maketitle
	% \maketitle also works in the article mode
	% \titlepage only works for presentation
\end{frame}

\begin{frame}[label=inhalt]{Overview}
	\tableofcontents
\end{frame}

\section{Introduction}

% only for the miniframe navigation
\subsection*{}

\begin{frame}
	\frametitle<presentation>{Introduction}
	\begin{block}{Why to give a presentation:}
		\begin{itemize}
			\item show the main arguments and results of your work
			\item produce interest to read the full paper/report
			\item goal: be educational and also entertaining
		\end{itemize}
	\end{block}
	\begin{block}{Advantages of using \LaTeX ~with the beamer package:}
		\begin{itemize}
			\item very easy if the report is already written in \LaTeX
			\item different themes which are usable in practice
			\item possibility to create handouts using \emph{beamerarticle}
		\end{itemize}
	\end{block}
\end{frame}

\section{Body}

\subsection{Equations}

\begin{frame}
	\frametitle<presentation>{Equations}
	\begin{block}{Pythagoras theorem:}
		\begin{equation}
			a^2 + b^2 = c^2 \label{eq:pythogoras}
		\end{equation}	\end{block}
	\begin{block}{It follows that:}
		\begin{align}
			a^2 &= c^2 - b^2 \\
			b^2 &= c^2 - a^2
		\end{align} 	
 	\end{block}
\end{frame}

\subsection{Figures}

\begin{frame}
	\frametitle<presentation>{Figures}
	\begin{figure}
		\centering
			\includegraphics[height=5cm]{unilogo}
		\caption{Logo of the university.}
		\label{fig:unilogo}
	\end{figure}
\end{frame}

\begin{frame}{Plots}
	\begin{figure}
		\centering
			\begin{tikzpicture}
				\begin{axis}[
					beamer,
					xlabel={time, $t$ (in \si{\milli\second})},
					ylabel={voltage, $u(t)$ (in \si{\volt})},
					xmin=0,xmax=20,
					ymin=-350,ymax=350,
					legend pos=south west,
					% example how to insert a line by giving a formula directly
					]
					\addplot+[
						domain=0:20,
						samples=101,
					] {sin(deg(x*2*pi/20))*sqrt(2)*230};
					\addlegendentry{sine wave-form};
					\addplot+[
						domain=0:20,
						samples=101,
					] {cos(deg(x*2*pi/20))*sqrt(2)*230};
					\addlegendentry{cosine wave-form};
				\end{axis}
			\end{tikzpicture}
		\caption{Harmonic time course of a voltage with a frequency of \SI{50}{\hertz} 
		and an effective value of \SI{230}{\volt}.}
		\label{fig:u_t_sinus}
	\end{figure}
\end{frame}

\subsection{Tables}

\begin{frame}
	\frametitle<presentation>{Tables}
	\begin{table}
		\centering
			\begin{tabular}{cc}
				\toprule
				variable & meaning \\
				\midrule
				$t$ & time \\
				$U$ & voltage \\
				\bottomrule
			\end{tabular}
		\caption{Example table.}
		\label{tab:example}
	\end{table}
\end{frame}

\begin{frame}{Citations}
	\begin{block}{Don't use short citations:}
		\begin{itemize}
			\item avoid short citations like~[1]
			\item no one will remember the numbers when the list of references is shown
			\item use full citations instead
		\end{itemize}
	\end{block}
	\begin{block}{Example for a full citation:}
		\small\fullcite{hering:technical_reports}
	\end{block}
\end{frame}

\section{Conclusion}

% only for the miniframe navigation
\subsection*{}

\begin{frame}
	\frametitle<presentation>{Conclusion}
	\begin{block}{Results:}
	  \begin{itemize}
	  	\item summarize the main results of your work
	  	\item also talk about remaining tasks or problems	  	
		\end{itemize}
	\end{block}
	\begin{block}{Questions:}
	 \begin{itemize}
	 	\item save some time for answering question
	 	\item optionally prepare some extra slides for supposable questions
	 \end{itemize}
	\end{block}	
\end{frame}

% only for the miniframe navigation
\section*{}

\begin{frame}<beamer>{}
	\begin{center}
		Thanks for your attention!
	\end{center}
	\begin{center}
		Are there questions?
	\end{center}
\end{frame}

\end{document}


\end{document}